\section{Introduction}\label{sec:intro}

The security of a computer system rests on its operating system. The operating
system ties software to hardware: it owns the devices, decides which process runs and when, and stands
between processes that must not learn about each other. Whatever an operating
system lets one process observe about another is a potential leak, and that
includes \emph{when} things happen, not only what values are computed.

A key part of this mediation is interrupt handling. Interrupts are how an operating
system learns that something happened: a disk transfer completed, a packet arrived,
a timer expired. Some of these events are secret. Whether a particular process just
received data from the disk may reveal what it is doing, so an operating system
that keeps such an event secret has to hide that it happened at all.

Ordinary interrupt handling does not hide it. As soon as the interrupt can be
serviced, its handler runs, displacing whatever process was scheduled; the handler
takes a few steps, and then the displaced process resumes. Suppose a public process keeps
issuing requests whose results an attacker can see. While the handler runs, those
requests stop, and the attacker sees a gap in the stream exactly as long as the
handler (\cref{fig:leak}). The interrupt has leaked through the \emph{schedule}.
Timing leaks through interrupt handling of this kind arise in several operating
systems \todo{examples, with citations}.

This is the example we follow throughout the paper. The secret is \emph{whether a
disk interrupt occurred}. A public process issues requests whose results the
attacker sees, and a private process, the legitimate consumer of the disk event,
reacts to it. The attacker may learn nothing about the disk interrupt. \Cref{sec:model,sec:designs,sec:proof}
model this system, exhibit the leak, and remove it.

\begin{figure}[t]
\centering
\begin{tabular}{lcccc}
\toprule
step & 1 & 2 & 3 & 4 \\
\midrule
no disk interrupt   & \code{Get} & \code{Get} & \code{Get} & \code{Get} \\
disk interrupt before step 1 & \code{Get} & ---   & ---   & \code{Get} \\
\bottomrule
\end{tabular}
\caption{What the attacker sees of the public process under ordinary interrupt
  handling, with handlers of two steps. The two runs differ only in whether a disk
  interrupt arrives, which the attacker may not learn. The interrupt sets a pending
  bit; step 1 was already scheduled, so the disk handler runs at steps 2 and 3, and
  the public process does not run then. The mechanised counterexample (\cref{thm:leak}) uses steps 1 and 2.}
\label{fig:leak}
\end{figure}

Leaks of this kind are the subject of information flow
control~\cite{sabelfeld2003}, where the absence of a leak is stated as
non-interference~\cite{goguen1982}: varying what the attacker may not see must not
change what it does see. For timing leaks, ``what it sees'' must include when each
output appears~\cite{agat2000}.

In this paper we use the process calculus of Rafnsson, Jia and
Bauer~\cite{rafnsson2017} to state and prove such properties of interrupt handling.
Its notion of non-interference is timing-sensitive, and its non-interference
theorems compose: if the parts of a process are non-interfering at suitable
equivalences, so is the whole. We model an operating system with interrupt handling
as one closed process in this calculus. Inside it, the two user processes (one
public, one private), the scheduler and one handler per kind of interrupt run as
concurrent processes in a pool, and a global state, holding the interrupt
controller, decides at each step which of them runs. Interrupts arrive as inputs to
the whole system, and at each step the pool emits the output of the process that
ran. We prove in
Rocq\footnote{The mechanisation is available at \url{\repobase}.} that ordinary
interrupt handling is interfering (\m{model_immediate_not_NI}): the leak of
\cref{fig:leak} is a counterexample.

We then model the corresponding non-interfering variant, obtained by altering how
interrupts are processed. The idea is to take away a secret interrupt's power to
choose \emph{when} its handler runs. Interrupts are serviced only inside time slices
opened by the public timer, every handler run takes the same number of steps, and
any part of a slice that no real interrupt claims is filled with a dummy handler,
the NOP handler. A secret handler run then replaces a NOP run instead of displacing a user
process, and the public schedule is the same whether or not the interrupt occurred.
We prove this variant non-interfering (\m{model_sliced_NI}), even to an attacker who
sees, at every step, which process in the pool ran, the scheduler and the handlers
included. Both variants are instances of a single generic model and differ only
in their \emph{interrupt policy}: the initial state and how handler completion is
serviced. The result is general in two ways. It holds for every scheduler
and every pair of user processes that are themselves non-interfering. And it holds
for every choice of the two numbers that fix the timing of the policy: the number
of steps one handler run takes, and the number of handler runs that fit in one
time slice.

The model fixes \emph{what} the modified interrupt processing must do; how to
implement it in an operating system is carried out in separate work
\todo{cite X}.

\paragraph{Contributions}
\begin{itemize}
\item A mechanised, timing-sensitive non-interference proof for interrupt handling
  (\cref{sec:designs,sec:proof}), and a mechanised counterexample for ordinary
  interrupt handling (\cref{sec:leak}). The proof is compositional: it applies one generic
  theorem per constructor of the model, and leaves one substantial obligation
  specific to interrupt handling (\cref{sec:fvni}).
\item Requirements, forced by the proof, on any implementation of interrupt
  handling that is to be non-interfering (\cref{sec:proof}). The arrival of a
  secret interrupt may only record that it is pending; every decision it influences, such as starting its handler, must be
  deferred to a point fixed by public time. Whether a secret handler ran must be
  indistinguishable from whether a dummy handler ran, so the time reserved for
  handlers must be fixed in advance and always used up. And the interrupt
  controller must be split by secrecy: which interrupts are pending is secret,
  while the masks that decide what runs next may depend only on public
  information. The implementation in \todo{cite X} follows these requirements.
\item A mechanisation of the calculus and its composition theorems. Our
  definitions depart from the original in four places, and for each we prove that
  the departure changes nothing or only weakens a side condition
  (\cref{sec:mechanisation}).
\end{itemize}

\Cref{sec:background} recalls the calculus and the definition of non-interference.
\Cref{sec:model} presents the generic model and \cref{sec:designs} the two interrupt policies.
\Cref{sec:proof} gives the proof, \cref{sec:generality} states precisely how general
it is, and \cref{sec:mechanisation} describes the mechanisation.
\Cref{sec:related,sec:conclusion} discuss related work and limitations.
